\subsection{Sprint 2}

Após uma boa entrega aos stakeholders, realizamos uma retrospectiva eficaz, onde foi
decidido mudar a metodologia de divisão das tarefas. Passamos a testar a entrega de
tarefas por pessoa: cada membro poderia visualizar as \textit{issues} e assinar a de
sua escolha. Assinei a tarefa de \textbf{criar professores}, baseando-me nas demais
entidades já criadas e completas, pois não tinha muita familiaridade com Spring Boot
--- essa seria minha primeira tarefa individual. Ao finalizar e entregar, a task não
passou na verificação por erros na lógica de negócio que poderiam impactar outras
entidades do projeto.

Além disso, o time adotou o \textbf{Swagger}, que trouxe muito mais consistência ao
código e ajudou na automação da inserção no banco de dados, permitindo entregas de
dados com menos erros nos \textit{endpoints}. Como o código precisava ser entregue
com urgência, meu líder técnico corrigiu os erros identificados.

Com todas as tarefas distribuídas, acompanhei algumas e decidi aprender a trabalhar no
front-end de forma independente. Desenvolvi duas telas: \textbf{planos de ensino} e o
\textbf{modal de lições}. Foi minha primeira experiência sem ajuda ou supervisão de
colegas. Eu sabia que não seria código inserido diretamente no projeto, pois um colega
do squad já havia assinado a issue --- foi um exercício de aprendizado pessoal.

Clonei o repositório, criei uma branch e comecei a desenvolver. Dediquei cerca de
\textbf{10 horas} para criar minha primeira tela \ac{html} do zero, deixá-la esteticamente
agradável, padronizada e com os containers corretamente posicionados. Um dos maiores
desafios foi entender como declarar a tela para que ela aparecesse na home da aplicação
--- a chamada não funcionava, e só depois de muita investigação identifiquei um erro
simples de digitação. Com a tela concluída, criei o modal de plano de ensino, que foi
mais fluido, pois eu já entendia o fluxo de criação.

Apesar de nenhum código meu ter sido aceito nesta sprint, ela representou um grande
aprendizado. Registro aqui também uma reflexão sobre o lançamento de horas: acredito
que ele deva refletir o serviço efetivamente realizado, com no máximo um estudo
pontual na Sprint 0. Não contabilizei todas as horas de estudo individual, pois
considero que o aprendizado faz parte da formação, não apenas da entrega.

Por experiência com programação ser quase exclusivamente acadêmica até então, há
momentos em que me sinto distante do nível técnico de alguns colegas --- muitos deles
são desenvolvedores plenos. No entanto, trabalhar com pessoas tão capacitadas é
exatamente o que me motiva a me esforçar além do normal.
